% Generated by scripts/assemble.py; corrections belong in source/corrections.json.
\examhead{171}{西南交通大学 2025 年数学分析试卷}{162}
% p163-b15; source PDF page 163
\blockanchor{p163-b15}
\question{1}(10 分) 证明: $\lim_{n \to \infty} \sin n$ 不存在.\par
% p163-b16; source PDF page 163
\blockanchor{p163-b16}
\question{2}(10 分) 设 $f(x)$ 是 $\mathbb{R}$ 上周期为 $T$ 的连续函数. 证明: $\lim_{t \to +\infty} \frac{1}{t} \int_{0}^{t} f(x) \, \mathrm{d}x = \frac{1}{T} \int_{0}^{T} f(x) \, \mathrm{d}x$.\par
% p163-b17; source PDF page 163
\blockanchor{p163-b17}
\question{3}(10 分) 设函数 $f(x)$ 在 $[1, +\infty)$ 上连续可微, 且 $\int_{1}^{+\infty} f(x) \mathrm{d}x$ 与 $\int_{1}^{+\infty} f'(x) \mathrm{d}x$ 均收敛, 证明: $\lim_{x \to +\infty} f(x) = 0$.\par
% p163-b18; source PDF page 163
\blockanchor{p163-b18}
\question{4}(15 分) 设函数 $f(x)$ 在 $(a, b)$ 内处处有导数 $f'(x)$. 证明: $(a, b)$ 中的点或者为 $f'(x)$ 的连续点, 或者为 $f'(x)$ 的第二类间断点.\par
% p163-b19; source PDF page 163
\blockanchor{p163-b19}
\question{5}(15 分) 设函数 $f(x)$ 在 $[1, +\infty)$ 上连续, 在 $(1, +\infty)$ 上可导, 已知函数 $e^{-x}f'(x)$ 在 $(1, +\infty)$ 上有界, 证明: 函数 $e^{-x}f(x)$ 在 $(1, +\infty)$ 上也有界.\par
% p163-b20; source PDF page 163
\blockanchor{p163-b20}
\question{6}(10 分) 求极限 $\lim_{x\to 0}\frac{\ln(1 + \sin^2 x) - 6(\sqrt[3]{2 - \cos x} - 1)}{x^4}$.\par
% p163-b21; source PDF page 163
\blockanchor{p163-b21}
\question{7}(10 分) 证明: 有限闭区间上的单调函数必定可积.\par
% p163-b22; source PDF page 163
\blockanchor{p163-b22}
\question{8}(10 分) 求幂级数 $\sum_{n=0}^{\infty} \frac{n^2 + 1}{2^n n!} x^n$ 的和函数.\par
% p164-b00; source PDF page 164
\blockanchor{p164-b00}
\question{9}(15 分) $K$ 为 $\mathbb{R}^n$ 中的紧集, $f: K \to \mathbb{R}^n$ 为连续映射, 证明: $f$ 在 $K$ 上一致连续.\par
% p164-b01; source PDF page 164
\par\Needspace{6\baselineskip}
\blockanchor{p164-b01}
\question{10}(15 分) 设 $u(x, y)$ 所有二阶偏导数都连续, 且\par
\nopagebreak[4]
% p164-b02; source PDF page 164
\blockanchor{p164-b02}
\displayformula{\frac {\partial^ {2} u}{\partial x ^ {2}} - \frac {\partial^ {2} u}{\partial y ^ {2}} = 0, u (x, 2 x) = x, u _ {x} (x, 2 x) = x ^ {2}.}
% p164-b03; source PDF page 164
\blockanchor{p164-b03}
求 $u_{xx}(x,2x), u_{xy}(x,2x), u_{yy}(x,2x)$.\par
% p164-b04; source PDF page 164
\blockanchor{p164-b04}
\question{11}(15 分) 计算积分 $I = \iiint_{V} \frac{\mathrm{d}V}{\rho^2}$, $\rho$ 为点 $(x, y, z)$ 到 $x$ 轴的距离, $V$ 表示一个棱台, 它的顶点坐标为 $A(0,0,1), B(0,1,1), C(1,1,1), D(0,0,2), E(0,2,2), F(2,2,2)$.\par
% p164-b05; source PDF page 164
\blockanchor{p164-b05}
\question{12}(15 分) 求 $\iint_{\Sigma} x \, \mathrm{d}y \, \mathrm{d}z + y \, \mathrm{d}z \, \mathrm{d}x + z \, \mathrm{d}x \, \mathrm{d}y,$$\Sigma$ 为上半球面 $z = \sqrt{R^2 - x^2 - y^2}, R > 0,$ 方向取上侧.\par
